\paragraph{E1: baseline reproduction.}
On OGW we reproduce the classic OBS+BSS + damage-index pipeline and the
feature-ML pipeline of \cite{schnur2022} to establish a trustworthy
anchor (target: near their reported leave-group-out accuracy).

\paragraph{E2: level-A eventization Pareto.}
For each frequency and damage position we sweep the SoD threshold
$\delta$ and record detection ROC/AUC against the event rate
(events per raw sample), producing the main Pareto front. At each rate
we compare against (i) uniform decimation and (ii) SVD/PCA compression
\cite{yang2023compression} at equal effective rate.

\paragraph{E3: temperature--damage discrimination.}
Without explicit compensation, the discriminator of
Section~\ref{sec:disc} separates temperature- from damage-induced
events; we additionally evaluate a temperature-extrapolation protocol
(baselines from 20--$40\,^\circ$C, test 40--$60\,^\circ$C) where
compensation-based pipelines degrade.

\paragraph{E4: long-term level-B validation.}
On the 4.5-year outdoor dataset we eventize the per-path DI series
(level B) and measure detection latency and false-alarm rate for the
three damage onsets (mass, dent, through-hole), reporting confounded
segments separately.

\paragraph{E5: directional analysis (CFRP).}
The 66 paths are binned by propagation direction; we assess whether
direction-dependent event-cluster strength/arrival carries usable
information for coarse localization, stated conservatively given the
quasi-isotropic layup.

\paragraph{E6: ablation.}
We ablate $\delta$, hysteresis, frequency selection, and compensation
quality (BSS vs.\ OBS vs.\ none) to map the sensitivity chain of the
eventization pipeline.
